\chapter{Wprowadzenie}
\label{chap:wprowadzenie}

Wprowadzenie osadza temat pracy w~szerszym kontekście: wyjaśnia, dlaczego
podjęty problem jest istotny z~punktu widzenia inżynierii środowiska lub
energetyki, jakie są obecne tendencje i~wymagania prawne oraz jakie
rozwiązania są stosowane w~praktyce. Przykładowo, w~pracy dotyczącej
ciepłownictwa można przywołać koncepcję sieci czwartej generacji
\cite{lund2014} oraz przegląd stanu ciepłownictwa na świecie
\cite{werner2017}, a~w~pracy o~budynkach -- wymagania dyrektywy
w~sprawie charakterystyki energetycznej budynków \cite{epbd2024}.

Poniżej zebrano wymagania edytorskie, które formatka spełnia automatycznie.
Przy pisaniu pracy nie trzeba ich ustawiać ręcznie -- wystarczy korzystać
z~poleceń opisanych w~rozdziale~\ref{chap:redakcja}.

\section{Wymagania edytorskie}
\label{sec:wymagania}

Podstawowe parametry składu zestawiono w~tabeli~\ref{tab:parametry}.
Każdy rozdział główny rozpoczyna się od nowej strony. Strona tytułowa jest
wliczana do numeracji, lecz numer nie jest na niej drukowany.

\begin{table}[htbp]
  \centering
  \caption{Parametry formatowania pracy dyplomowej}
  \label{tab:parametry}
  \begin{tabularx}{\linewidth}{@{}>{\raggedright\arraybackslash}p{4.4cm}X@{}}
    \toprule
    \textbf{Element} & \textbf{Ustawienie} \\
    \midrule
    Format strony      & A4, druk jedno- lub dwustronny \\
    Marginesy          & lewy (wewnętrzny) \qty{3.5}{cm}, pozostałe \qty{2.5}{cm} \\
    Krój i stopień     & Krój typu Times lub Arial, \qty{11}{pt} \\
    Interlinia         & 1,25 wiersza (tekst), pojedyncza (podpisy, tabele, kod) \\
    Akapit             & wyjustowany, wcięcie pierwszego wiersza \qty{0.63}{cm} \\
    Tytuły rozdziałów  & \qty{14}{pt}, pogrubione, numeracja „1.” \\
    Podrozdziały       & \qty{13}{pt} i~\qty{12}{pt}, pogrubione, numeracja „1.1.” i~„1.1.1.” \\
    Podpisy            & kursywa, rysunek -- pod obiektem, tabela -- nad obiektem \\
    Numeracja obiektów & w~obrębie rozdziału: rys.~3.1, tabela~3.1, wzór~(3.1) \\
    Numer strony       & na dole strony, wyśrodkowany \\
    \bottomrule
  \end{tabularx}
\end{table}

\section{Druk i oprawa}

Jeżeli jeden z~egzemplarzy ma być wydrukowany dwustronnie (np.\ w~miękkiej
oprawie), należy w~pliku \texttt{praca.tex} zmienić opcję klasy
z~\texttt{jednostronnie} na \texttt{dwustronnie}. Marginesy staną się
wówczas lustrzane (szerszy margines zawsze od strony grzbietu), rozdziały
będą zaczynały się na stronach nieparzystych, a~wstawione strony puste nie
otrzymają numeru. Wersję do archiwizacji w~systemie ASAP zwykle składa się
jednostronnie -- warto potwierdzić to z~promotorem.
